\documentclass[aspectratio=169,11pt]{beamer}
\usepackage[utf8]{inputenc}
\usepackage[T1]{fontenc}
\usepackage{lmodern}
\usepackage[english]{babel}
\usepackage{graphicx}
\usepackage{booktabs}
\usepackage{amsmath,amssymb}
\usepackage{xcolor}
\usepackage{listings}
\usetheme{Madrid}
\usecolortheme{default}
\setbeamertemplate{navigation symbols}{}
\setbeamertemplate{footline}[frame number]
\setbeamertemplate{blocks}[rounded][shadow=false]
\AtBeginSection[]{\begin{frame}{Outline}\tableofcontents[currentsection,subsectionstyle=hide]\end{frame}}
\definecolor{hsbiblue}{RGB}{45,57,120}
\definecolor{hsbimid}{RGB}{76,114,176}
\definecolor{hsbilight}{RGB}{225,231,247}
\definecolor{cgreen}{RGB}{85,164,103}
\definecolor{cred}{RGB}{196,78,82}
\lstset{basicstyle=\ttfamily\scriptsize,breaklines=true,keywordstyle=\color{hsbiblue},commentstyle=\color{cgreen},columns=fullflexible,showstringspaces=false}
\title{Fast Track --- Mock Exam 1 (Solutions)}
\subtitle{CloudDesk scenario $\cdot$ full worked solutions}
\author{Prof.\ Dr.\ Christoph Weisser}
\institute{HSBI}
\date{Summer Semester 2026}
\begin{document}
\begin{frame}\titlepage\end{frame}

\begin{frame}{How to use these solutions}
\begin{itemize}
  \item Companion to \texttt{exam\_1.md} --- the CloudDesk parallel-form paper (100 marks, 120 min).
  \item \textbf{Five parts, 20 marks each}; one slide per question. The correct MCQ option is stated
        \emph{and} every distractor is explained.
  \item \textbf{Green} = the key point / correct answer; \textbf{red} = the trap.
  \item Code shown is one correct solution --- equivalent idioms earn full marks.
  \item Model marks are per sub-part; method marks apply on calculations.
\end{itemize}
\vspace{0.4em}
\begin{block}{Grade banding (guide)}
\footnotesize $\geq 85$ distinction \; $\cdot$ \; $70$--$84$ strong \; $\cdot$ \; $50$--$69$ pass \; $\cdot$ \; $<50$ revisit the flagged notebooks.
\end{block}
\end{frame}

% ==================================================================
\section{Part A --- Python foundations}
% ==================================================================

\begin{frame}[fragile]{A1 --- \texttt{int()} truncates toward zero (2 marks)}
\textbf{Q.} What does \texttt{print(int(-3.9))} output? \hfill \textbf{Answer: B) \texttt{-3}}
\begin{itemize}
  \item \texttt{int()} on a float \textcolor{cred}{truncates toward zero} --- it drops the fractional
        part, it does \emph{not} round. So \texttt{-3.9} $\to$ \texttt{-3}.
  \item \textbf{A) \texttt{-4}} --- that would be \emph{rounding} (or \texttt{math.floor}); floor of
        $-3.9$ is $-4$, but \texttt{int} does not floor negatives.
  \item \textbf{C) \texttt{-3.9}} --- \texttt{int()} always returns an \texttt{int}, never a float.
  \item \textbf{D) ValueError} --- only strings that aren't numeric raise; a float is always castable.
\end{itemize}
\begin{block}{Remember}
\footnotesize \texttt{int(3.9) == 3} and \texttt{int(-3.9) == -3}. Use \texttt{round()} to round to nearest.
\end{block}
\end{frame}

\begin{frame}[fragile]{A2 --- f-string percent format (2 marks)}
\textbf{Q.} \texttt{rate = 0.234; f"\{rate:.1\%\}"} \hfill \textbf{Answer: B) \texttt{"23.4\%"}}
\begin{itemize}
  \item The \texttt{\%} format spec \textcolor{cgreen}{multiplies by 100 and appends \texttt{\%}};
        \texttt{.1} keeps one decimal place: $0.234 \to 23.4\%$.
  \item \textbf{A) \texttt{"0.2\%"}} --- that ignores the $\times 100$ and just does \texttt{.1f} with a
        percent sign; wrong.
  \item \textbf{C) \texttt{"23\%"}} --- that is \texttt{:.0\%} (zero decimals); the spec asked for one.
  \item \textbf{D) \texttt{"0.234\%"}} --- treats the number as already-a-percentage; \texttt{\%} always scales.
\end{itemize}
\begin{block}{Related specs}
\footnotesize \texttt{\{x:,\}} thousands separator, \texttt{\{x:.2f\}} fixed decimals, \texttt{\{x:>12\}} right-align.
\end{block}
\end{frame}

\begin{frame}[fragile]{A3 --- aliasing, \texttt{is} vs \texttt{==} (3 marks)}
\begin{lstlisting}[language=Python]
a = [1, 2, 3]
b = a            # b is a second NAME for the SAME list object
b.append(4)
print(len(a))    # -> 4
\end{lstlisting}
\begin{itemize}
  \item \textbf{Output: \texttt{4}.} \texttt{b = a} does not copy; both names point at one list, so
        mutating through \texttt{b} is visible through \texttt{a}. \hfill \emph{(1.5)}
  \item \textcolor{hsbiblue}{\texttt{==}} tests \emph{value} equality (same contents);
        \textcolor{hsbiblue}{\texttt{is}} tests \emph{identity} (same object in memory,
        same \texttt{id()}). Here \texttt{a is b} is \texttt{True}. \hfill \emph{(1.5)}
  \item To get an independent copy: \texttt{b = a.copy()} / \texttt{list(a)} / \texttt{a[:]}
        (and \texttt{copy.deepcopy} for nested lists).
\end{itemize}
\end{frame}

\begin{frame}[fragile]{A4 --- most common channel with \texttt{Counter} (4 marks)}
\begin{lstlisting}[language=Python]
from collections import Counter

counts = Counter(channels)          # {'chat': 12, 'phone': 7, ...}
channel, n = counts.most_common(1)[0]
# most_common(1) -> [('chat', 12)]; [0] -> ('chat', 12); unpack
\end{lstlisting}
\begin{itemize}
  \item \texttt{Counter(channels)} tallies each channel in one pass. \hfill \emph{(2)}
  \item \texttt{.most\_common(1)} returns a \emph{list} with the single top
        \texttt{(value, count)} tuple; index \texttt{[0]} then unpack. \hfill \emph{(2)}
  \item \textcolor{cgreen}{Accept} equivalents, e.g. \texttt{max(counts, key=counts.get)} for the
        channel plus \texttt{counts[channel]} for the count.
\end{itemize}
\end{frame}

\begin{frame}[fragile]{A5 --- mutable default argument (4 marks)}
\textbf{Bug:} the default \texttt{batch=[]} is created \textcolor{cred}{once, when the function is
defined}, and \emph{shared across all calls} --- so appends accumulate. \hfill \emph{(2)}
\begin{columns}[T]
\column{0.48\textwidth}
\begin{block}{Broken}
\begin{lstlisting}[language=Python]
def log_ticket(ticket_id, batch=[]):
    batch.append(ticket_id)
    return batch
\end{lstlisting}
\end{block}
\column{0.48\textwidth}
\begin{block}{Fixed}
\begin{lstlisting}[language=Python]
def log_ticket(ticket_id, batch=None):
    if batch is None:
        batch = []
    batch.append(ticket_id)
    return batch
\end{lstlisting}
\end{block}
\end{columns}
\vspace{0.4em}
Use \texttt{None} as the sentinel and build a fresh list inside the body. \hfill \emph{(2)}
\end{frame}

\begin{frame}[fragile]{A6 --- \texttt{SupportTicket} class (5 marks)}
\begin{lstlisting}[language=Python]
class SupportTicket:
    count = 0                          # class attribute (shared) -- (1)

    def __init__(self, ticket_id, csat):
        self.ticket_id = ticket_id     # instance attributes  -- (2)
        self.csat = csat
        SupportTicket.count += 1       # bump the shared counter -- (1)

    def is_happy(self):
        return self.csat >= 4          # method -- (1)
\end{lstlisting}
\begin{itemize}
  \item \textcolor{cred}{Common slip:} \texttt{self.count += 1} creates an \emph{instance}
        attribute and leaves the class counter at 0 --- increment \texttt{SupportTicket.count}
        (or \texttt{type(self).count}).
  \item \texttt{SupportTicket("T1", 5).is\_happy()} $\to$ \texttt{True}; \texttt{SupportTicket.count} grows by 1 per instance.
\end{itemize}
\end{frame}

% ==================================================================
\section{Part B --- Data, pandas, viz \& statistics}
% ==================================================================

\begin{frame}[fragile]{B1 --- boolean masks need parentheses (2 marks)}
\textbf{Answer: B)} \texttt{df[(df["csat"] < 3) \& (df["channel"] == "phone")]}
\begin{itemize}
  \item Combine masks with the \emph{element-wise} operator \textcolor{hsbiblue}{\texttt{\&}}, and
        \textcolor{cgreen}{wrap each condition in parentheses} --- \texttt{\&} binds tighter than
        \texttt{<}/\texttt{==}.
  \item \textbf{A)} uses Python \texttt{and} --- raises \emph{"truth value of a Series is ambiguous"}.
  \item \textbf{C)} no parentheses: \texttt{3 \& df[...]} is evaluated first $\Rightarrow$ wrong /
        \texttt{TypeError}.
  \item \textbf{D)} the comma makes pandas read it as \texttt{(rows, cols)} indexing --- not a mask.
\end{itemize}
\end{frame}

\begin{frame}{B2 --- seaborn error bars (2 marks)}
\textbf{Answer: B) the 95\% confidence interval of the mean.}
\begin{itemize}
  \item By default \texttt{sns.barplot} plots the \emph{mean} of \texttt{y} per \texttt{x} category,
        with error bars = the \textcolor{cgreen}{95\% CI} (bootstrapped) of that mean.
  \item \textbf{A) $\pm$1 std} --- available via \texttt{errorbar=("sd", 1)}, but not the default.
  \item \textbf{C) min--max range} --- seaborn never uses full range by default.
  \item \textbf{D) IQR} --- that describes a box plot, not a bar plot's error bar.
\end{itemize}
\begin{block}{Point}
\footnotesize The bar shows an \emph{estimate of the mean} with its uncertainty --- overlapping CIs hint the difference may not be significant.
\end{block}
\end{frame}

\begin{frame}[fragile]{B3 --- mean of a 0/1 column is a rate (3 marks)}
\texttt{tickets.groupby("channel")["escalated"].mean()}
\begin{itemize}
  \item Returns a \textbf{Series indexed by channel}, each value = the \emph{mean of
        \texttt{escalated} within that channel}. \hfill \emph{(1.5)}
  \item For a 0/1 column, $\text{mean} = \dfrac{\text{number of 1s}}{\text{number of rows}}
        = \dfrac{\#\text{escalated}}{\#\text{tickets}}$ --- exactly the \textcolor{cgreen}{escalation
        rate} (a proportion in $[0,1]$). \hfill \emph{(1.5)}
  \item So \texttt{0.18} for \texttt{phone} means 18\% of phone tickets escalated. (This is the same
        trick used to read channel/plan escalation rates in the capstone.)
\end{itemize}
\end{frame}

\begin{frame}[fragile]{B4 --- mean CSAT per channel, sorted (4 marks)}
\begin{lstlisting}[language=Python]
tickets.groupby("channel")["csat"].mean().sort_values()
\end{lstlisting}
\begin{itemize}
  \item \texttt{groupby("channel")["csat"].mean()} $\to$ one mean per channel. \hfill \emph{(2)}
  \item \texttt{.sort\_values()} sorts \emph{ascending} by default, so the worst channel is
        first. \hfill \emph{(2)}
  \item \textcolor{cred}{Trap:} \texttt{.sort\_values(ascending=False)} would put the \emph{best}
        first --- the question asked lowest-first.
\end{itemize}
\end{frame}

\begin{frame}[fragile]{B5 --- SQL: GROUP BY + HAVING (4 marks)}
\begin{lstlisting}[language=SQL]
SELECT channel, AVG(csat) AS avg_csat
FROM tickets
GROUP BY channel
HAVING COUNT(*) > 100
ORDER BY avg_csat DESC;
\end{lstlisting}
\begin{itemize}
  \item \texttt{GROUP BY channel} then \texttt{AVG(csat)} per group. \hfill \emph{(1.5)}
  \item \textcolor{cgreen}{\texttt{HAVING}} filters \emph{groups} by an aggregate
        (\texttt{COUNT(*) > 100}) --- \texttt{WHERE} cannot, it filters raw rows before
        grouping. \hfill \emph{(1.5)}
  \item \texttt{ORDER BY avg\_csat DESC} sorts highest-first. \hfill \emph{(1)}
\end{itemize}
\end{frame}

\begin{frame}{B6 --- Welch t-test interpretation (5 marks)}
\textbf{(a)} It is a \textbf{two-sample, two-sided Welch's t-test}.
\texttt{equal\_var=False} means it \emph{does not assume the two groups have equal variance}
(it uses the Welch--Satterthwaite d.f.) --- the safer default when group sizes/spreads
differ. \hfill \emph{(2)}

\textbf{(b)} $p = 0.004 < 0.05$, so we \textcolor{cgreen}{reject $H_0$}: the mean CSAT of chat and
phone \emph{differ significantly}. The negative statistic says chat's mean is the lower of the
two. \hfill \emph{(2)}

\textbf{(c)} Statistical $\neq$ practical significance: the \textbf{effect size} (e.g. Cohen's d)
could be tiny, so the CSAT gap may be too small to justify any operational change / cost.
\hfill \emph{(1)}
\end{frame}

% ==================================================================
\section{Part C --- ML, evaluation \& feature engineering}
% ==================================================================

\begin{frame}[fragile]{C1 --- what is data leakage (2 marks)}
\textbf{Answer: C) fitting \texttt{StandardScaler} on the whole dataset before splitting.}
\begin{itemize}
  \item That lets the scaler's \texttt{mean\_}/\texttt{scale\_} \textcolor{cred}{peek at the test
        rows} --- test information leaks into training. Fit preprocessing on \emph{train only}
        (ideally inside a \texttt{Pipeline}).
  \item \textbf{A)} fitting on train only is exactly right --- no leak.
  \item \textbf{B)} \texttt{stratify=y} just preserves class balance across the split --- good practice.
  \item \textbf{D)} a \texttt{Pipeline} \emph{prevents} leakage by re-fitting the scaler per CV fold.
\end{itemize}
\end{frame}

\begin{frame}[fragile]{C2 --- \texttt{confusion\_matrix().ravel()} order (2 marks)}
\textbf{Answer: B) \texttt{tn, fp, fn, tp}.}
\begin{itemize}
  \item For labels \texttt{[0, 1]} the matrix is row = \emph{true}, column = \emph{predicted};
        flattening row-major gives \textcolor{cgreen}{\texttt{tn, fp, fn, tp}}.
  \item \textbf{A)} reversed order (tp first) --- wrong; that is not how \texttt{.ravel()} lays it out.
  \item \textbf{C)} swaps \texttt{fp} and \texttt{fn} --- a classic mix-up that flips your cost sums.
  \item \textbf{D)} those are \texttt{classification\_report} columns, not confusion-matrix cells.
\end{itemize}
\begin{block}{Recall}
\footnotesize FN = a real churner/fraud you \emph{missed}; FP = a false alarm. The costs usually differ a lot.
\end{block}
\end{frame}

\begin{frame}{C3 --- why 99.5\% accuracy is misleading (3 marks)}
\begin{itemize}
  \item The classes are \textbf{severely imbalanced} ($\sim$0.5\% fraud). A model that always
        predicts "not fraud" is right 99.5\% of the time \textcolor{cred}{yet catches zero
        fraud} --- accuracy rewards the majority class and hides total failure on the class we
        care about. \hfill \emph{(2)}
  \item Better metrics that expose it: \textbf{recall} (fraction of fraud actually caught, here
        $0$), \textbf{PR-AUC / average precision}, or a \textbf{cost-weighted} confusion matrix.
        \hfill \emph{(1)}
  \item Any \emph{one} named correctly earns the mark; ROC-AUC is acceptable but PR-AUC is
        preferred when positives are rare.
\end{itemize}
\end{frame}

\begin{frame}[fragile]{C4 --- standardise + logistic regression pipeline (4 marks)}
\begin{lstlisting}[language=Python]
from sklearn.pipeline import make_pipeline
from sklearn.preprocessing import StandardScaler
from sklearn.linear_model import LogisticRegression

model = make_pipeline(StandardScaler(),
                      LogisticRegression(max_iter=1000))
model.fit(X_train, y_train)
proba = model.predict_proba(X_test)[:, 1]   # P(churn) = positive class
\end{lstlisting}
\begin{itemize}
  \item Pipeline chains scaler $\to$ model, so the scaler is fit on train only. \hfill \emph{(2)}
  \item \texttt{predict\_proba(...)[:, 1]} takes column 1 = probability of the \textcolor{cgreen}{positive}
        (churn) class --- \emph{not} \texttt{predict}, which returns hard 0/1 labels. \hfill \emph{(2)}
\end{itemize}
\end{frame}

\begin{frame}{C5 --- expected-value offer threshold (5 marks)}
\textbf{(a) CLV} $= m / c = 300 / 0.02 = \textcolor{hsbiblue}{\text{EUR }15{,}000}$. \hfill \emph{(1)}

\textbf{(b) Break-even probability}
\[
p^{*} = \frac{C_{\text{offer}}}{s \cdot \text{CLV}}
      = \frac{60}{0.30 \times 15{,}000}
      = \frac{60}{4{,}500}
      = \textcolor{hsbiblue}{0.0133}\;(\approx 1.3\%). \hfill (2)
\]

\textbf{(c)} Predicted churn $0.06 > p^{*}=0.0133$, so the offer's expected value is
\emph{positive} --- \textcolor{cgreen}{send it}. Check:
\[
EV = p \cdot s \cdot \text{CLV} - C_{\text{offer}}
   = 0.06 \times 0.30 \times 15{,}000 - 60 = 270 - 60 = \text{EUR }210 > 0. \hfill (2)
\]
\begin{block}{Key idea}
\footnotesize The threshold is \emph{per-customer} --- a flat 0.5 cutoff is wrong; high-CLV customers justify offers at very low churn risk.
\end{block}
\end{frame}

\begin{frame}[fragile]{C6 --- target leakage via group statistic (4 marks)}
\textbf{Bug:} \texttt{plan\_churn\_rate} is computed with \texttt{transform("mean")} over
\textcolor{cred}{all rows (train + test) using the target \texttt{churned}} --- so each row's
feature already encodes label information from the test set. Leakage. \hfill \emph{(2)}
\begin{lstlisting}[language=Python]
# Fix: split FIRST, learn per-plan means on TRAIN only, map onto test.
X = data.drop(columns=["churned"]); y = data["churned"]
X_tr, X_te, y_tr, y_te = train_test_split(X, y, test_size=0.2,
                                          random_state=42)
rates = y_tr.groupby(X_tr["plan"]).mean()          # train-only means
X_tr["plan_churn_rate"] = X_tr["plan"].map(rates)
X_te["plan_churn_rate"] = X_te["plan"].map(rates)  # reuse train stats
\end{lstlisting}
\footnotesize Equivalent: a per-fold \texttt{TargetEncoder} inside a \texttt{Pipeline}. \hfill \emph{(2)}
\end{frame}

% ==================================================================
\section{Part D --- AI engineering}
% ==================================================================

\begin{frame}{D1 --- MCP primitive controlled by the model (2 marks)}
\textbf{Answer: C) tools.}
\begin{itemize}
  \item MCP's three primitives map to \emph{who is in control}: \textcolor{cgreen}{\textbf{tools} ---
        model-controlled} (the model decides to call them, like a POST endpoint with side effects).
  \item \textbf{A) resources} --- \emph{application}-controlled, read-only data (like a GET / file).
  \item \textbf{B) prompts} --- \emph{user}-controlled, invoked deliberately (like a slash command).
  \item \textbf{D) sampling} --- not one of the three core primitives (it is a server$\to$host request
        for a completion).
\end{itemize}
\begin{block}{Mnemonic}
\footnotesize tools = model, resources = app, prompts = user.
\end{block}
\end{frame}

\begin{frame}{D2 --- cosine of normalised vectors (2 marks)}
\textbf{Answer: B) their dot product.}
\begin{itemize}
  \item $\cos(a,b)=\dfrac{a\cdot b}{\lVert a\rVert\,\lVert b\rVert}$. If both are
        \textbf{L2-normalised} then $\lVert a\rVert=\lVert b\rVert=1$, so
        $\cos(a,b)=\textcolor{cgreen}{a\cdot b}$. This is why normalised retrieval is just
        \texttt{doc\_emb @ q}.
  \item \textbf{A) Euclidean distance} --- related but not equal; distance grows as similarity falls.
  \item \textbf{C) always 1} --- only if the vectors are identical in direction.
  \item \textbf{D) sum of elements} --- unrelated to the dot product between two vectors.
\end{itemize}
\end{frame}

\begin{frame}{D3 --- parse defensively; shape vs truth (3 marks)}
\textbf{(a)} An LLM can return prose, markdown fences, or malformed JSON, so \texttt{json.loads}
may raise \texttt{JSONDecodeError}. Wrapping it in \texttt{try/except} lets the pipeline
\textcolor{cgreen}{fall back} (e.g. \texttt{\{"topic": "unknown"\}} or a retry) instead of
crashing on one bad reply. \hfill \emph{(1.5)}

\textbf{(b)} Schema validation only checks the \emph{form}: correct keys, types, ranges,
non-empty. \texttt{\{"order\_id": "INV", "amount": 42.0\}} is well-\emph{shaped} --- a non-empty
string and a positive number --- so it \textcolor{cgreen}{passes}, even though the id is truncated
and \emph{factually wrong}. Only comparison with \textbf{ground truth} (a golden set) catches
wrongness. \emph{"Validation checks shape, not truth."} \hfill \emph{(1.5)}
\end{frame}

\begin{frame}[fragile]{D4 --- top-k by cosine similarity (4 marks)}
\begin{lstlisting}[language=Python]
import numpy as np

def top_k(query_vec, doc_matrix, k=3):
    q = query_vec / np.linalg.norm(query_vec)
    d = doc_matrix / np.linalg.norm(doc_matrix, axis=1, keepdims=True)
    sims = d @ q                      # cosine sim per row
    return np.argsort(-sims)[:k]      # indices, highest first
\end{lstlisting}
\begin{itemize}
  \item Normalise the query and \emph{each row} (\texttt{axis=1, keepdims=True}) so the dot product
        equals cosine similarity. \hfill \emph{(2)}
  \item \texttt{np.argsort(-sims)} sorts \emph{descending}; slice \texttt{[:k]}. \hfill \emph{(2)}
  \item \textcolor{cgreen}{Also accepted:} \texttt{sklearn.metrics.pairwise.cosine\_similarity}.
\end{itemize}
\end{frame}

\begin{frame}{D5 --- Mean Reciprocal Rank (3 marks)}
\textbf{(a)} Reciprocal rank of the first relevant hit per query:
\[
\text{RR}_1 = \tfrac{1}{1}=1,\quad \text{RR}_2 = \tfrac{1}{3}\approx 0.333,\quad \text{RR}_3 = 0 \ (\text{no hit}).
\]
\[
\text{MRR} = \frac{1 + 0.333 + 0}{3} = \frac{1.333}{3} \approx \textcolor{hsbiblue}{0.444}. \hfill (2)
\]

\textbf{(b)} MRR captures \emph{how high} the first correct document is ranked (rank 1 beats rank 3),
whereas a plain \texttt{retrieval@k} hit-rate only asks \emph{whether} a correct doc is somewhere in
the top-$k$, ignoring its position. \hfill \emph{(1)}
\end{frame}

\begin{frame}{D6 --- LLM-as-judge biases (6 marks)}
\textbf{(a) Two biases} (any two; 1.5 each): \hfill \emph{(3)}
\begin{itemize}
  \item \textbf{Verbosity / length bias} --- the judge favours longer, more confident-sounding
        answers even when they are no more correct.
  \item \textbf{Position / order bias} --- when comparing two answers, it tends to prefer whichever
        is shown \emph{first}.
\end{itemize}
\textbf{(b) Mitigations} (match to the bias; 1 each): \hfill \emph{(2)}
\begin{itemize}
  \item Verbosity $\to$ constrain the rubric / penalise length, or use a stronger/different judge model.
  \item Position $\to$ \textcolor{cgreen}{swap the order and average both runs}.
\end{itemize}
\textbf{(c)} \textbf{No} --- 62\% agreement is below the $\sim$80\% human-agreement bar, so the judge
is not yet reliable enough to gate deployments. \hfill \emph{(1)}
\end{frame}

% ==================================================================
\section{Part E --- Applied \& production}
% ==================================================================

\begin{frame}[fragile]{E1 --- cron: every 15 min on weekdays (2 marks)}
\textbf{Answer: C) \texttt{*/15 * * * 1-5}.}
\begin{itemize}
  \item Fields: \texttt{minute hour day-of-month month day-of-week}.
        \texttt{*/15} = every 15 min; \texttt{1-5} = Mon--Fri (cron uses
        \textcolor{cgreen}{Sunday = 0}, so Mon--Fri is \texttt{1-5}).
  \item \textbf{A) \texttt{0-4}} --- that is Sun--Thu (off by one); Friday is missed, Sunday wrongly included.
  \item \textbf{B) \texttt{15 * * * 1-5}} --- fires \emph{once} per hour at \texttt{:15}, not every 15 min.
  \item \textbf{D) \texttt{6-7}} --- Saturday (6) only (\texttt{7} is out of range / rarely Sunday); wrong days.
\end{itemize}
\end{frame}

\begin{frame}{E2 --- stop words drop "not" (2 marks)}
\textbf{Answer: B) \texttt{"not"} is an English stop word, so it is dropped before the model sees it.}
\begin{itemize}
  \item \texttt{TfidfVectorizer(stop\_words="english")} \textcolor{cred}{removes "not"}, so
        \emph{"not helpful"} and \emph{"helpful"} vectorise almost identically --- the negation is
        silently lost.
  \item \textbf{A)} false --- logistic regression handles TF-IDF text features fine (it is the workhorse rung).
  \item \textbf{C)} false --- TF-IDF output is L2-normalised by default.
  \item \textbf{D)} false --- length is not the issue; the phrase vectorises, just without "not".
\end{itemize}
\end{frame}

\begin{frame}{E3 --- Docker layer ordering (3 marks)}
\textbf{(a)} Docker caches layers top-to-bottom and \textcolor{cred}{invalidates every layer below
the first change}. Dependencies change \emph{rarely}, code changes \emph{constantly}. Putting
\texttt{RUN pip install} above \texttt{COPY src/} means editing code does \emph{not} invalidate the
(slow) install layer --- it stays cached, so rebuilds are fast. \hfill \emph{(2)}

\textbf{(b)} Editing only a source file rebuilds \textbf{$\sim$2 layers} (the \texttt{COPY src/}
and the \texttt{CMD} below it). Bumping \texttt{requirements.txt} invalidates from the
\texttt{COPY requirements.txt} line downward --- \textbf{$\sim$4 layers}, including the expensive
\texttt{pip install}. \hfill \emph{(1)}
\end{frame}

\begin{frame}[fragile]{E4 --- MAE and MAPE (4 marks)}
\begin{lstlisting}[language=Python]
import numpy as np
y_true = np.array([100, 120,  90, 110])
y_pred = np.array([110, 115, 100, 105])

mae  = np.abs(y_true - y_pred).mean()                       # 7.5
mape = (np.abs(y_true - y_pred) / np.abs(y_true)).mean() * 100  # ~7.46%
\end{lstlisting}
\begin{itemize}
  \item MAE = mean of $|y-\hat y|$: $(10+5+10+5)/4 = 7.5$. \hfill \emph{(2)}
  \item MAPE = mean of $|y-\hat y|/|y|$, $\times 100$:
        $(0.10+0.0417+0.1111+0.0455)/4\times100 \approx 7.46\%$. \hfill \emph{(2)}
  \item \textcolor{cgreen}{Note:} MAPE divides by the actuals, so it blows up if any actual is near 0
        --- prefer MAE when zeros are possible.
\end{itemize}
\end{frame}

\begin{frame}[fragile]{E5 --- \texttt{NoneType} has no \texttt{.get\_text} (4 marks)}
\textbf{Error:} \texttt{AttributeError: 'NoneType' object has no attribute 'get\_text'}.
\texttt{select\_one(".priority")} returns \textcolor{cred}{\texttt{None}} when a card has no such
element, and \texttt{None.get\_text(...)} then crashes. \hfill \emph{(2)}
\begin{lstlisting}[language=Python]
badge = soup.select_one(".priority")
priority = badge.get_text(strip=True) if badge else "normal"
\end{lstlisting}
\begin{itemize}
  \item Capture the node, \textcolor{cgreen}{test for \texttt{None}}, then read it. \hfill \emph{(2)}
  \item Defensive parsing turns a 3 a.m.\ crash into a sane default you can monitor.
\end{itemize}
\end{frame}

\begin{frame}{E6 --- costed decision rule (5 marks)}
\textbf{(a)} Total cost from the confusion matrix:
\[
\text{cost} = \text{FN}\cdot C_{\text{miss}} + \text{FP}\cdot C_{\text{false}}
            = 15\cdot 40 + 40\cdot 20 = 600 + 800 = \textcolor{hsbiblue}{\text{EUR }1{,}400}. \hfill (2)
\]

\textbf{(b)} \emph{Escalate nothing} $\Rightarrow$ every real escalation is an FN (huge miss cost);
\emph{escalate everything} $\Rightarrow$ a flood of FPs. Instead, \textcolor{cgreen}{sweep the
threshold} (\texttt{np.linspace}), price each confusion matrix with the formula above, and pick the
threshold with \texttt{np.argmin(costs)}. \hfill \emph{(2)}

\textbf{(c)} The headline is the \textbf{euro cost / saving of the rule} --- "the number that leaves
the room." AUC is a \emph{ranking-quality} metric, not a decision, so it can't be spent or
budgeted. \hfill \emph{(1)}
\end{frame}

\begin{frame}{Mark summary}
\begin{center}
\begin{tabular}{llc}
\toprule
\textbf{Part} & \textbf{Coverage} & \textbf{Marks} \\
\midrule
A & Python foundations (NB 1--5) & 20 \\
B & Data, pandas, viz \& statistics (NB 6--9) & 20 \\
C & ML, evaluation \& feature engineering (NB 8, 16, 17) & 20 \\
D & AI engineering: LLM, RAG, tools/agents, MCP, doc AI (NB 10--14, 18) & 20 \\
E & Applied \& production: time series, NLP, deploy, scraping, capstone & 20 \\
\midrule
\multicolumn{2}{l}{\textbf{Total} --- 30 questions} & \textbf{100} \\
\bottomrule
\end{tabular}
\end{center}
\vspace{0.8em}
\begin{center}\small End of solutions --- CloudDesk Mock Exam 1.\end{center}
\end{frame}

\end{document}
